\section{Methods}\label{sec6}

\subsection{Study design}

This study is an observational analysis of naturalistic public human--LLM conversations from WildChat-4.8M~\cite{zhao2024wildchat} (reported below as WildChat) , LMSYS Chat-1M~\cite{zheng2023lmsyschat1m} (reported below as LMSYS Chat) and the ShareChat~\cite{yan2025sharechat} strict-English subset (reported below as ShareChat) across two everyday task ecologies, coding and writing. It asks whether learning-oriented engagement is visible in ordinary use, where constructive engagement is concentrated and how assistant scaffolding is ordered with user engagement across conversation-level and adjacent-turn timescales. The design is descriptive and associational, following the logic of large-scale digital trace and computational social-science analyses~\cite{lazer2009computational,salganik2017bitbybit}. It does not estimate randomized exposure effects, causal instructional efficacy or durable learning outcomes~\cite{hernan2020causal}.

The empirical sequence mirrors the Results section. We first estimate the prevalence and composition of user engagement, then examine how constructive engagement varies by user framing, task ecology and interaction depth. We next test whether scaffolded assistant support marks conversations with richer constructive participation, separate scaffolded support by support form and analyse adjacent-turn coupling between assistant support and user engagement. Dataset and task ecology are treated as substantive dimensions of the interaction environment rather than nuisance variables to be averaged away.

\subsection{Data sources and analytic samples}

The analytic corpus contains 128,569 conversations and 981,470 total turns, including 491,685 user turns and 489,785 assistant turns. The six task settings are WildChat coding, WildChat writing, LMSYS Chat coding, LMSYS Chat writing, ShareChat coding and ShareChat writing. WildChat and LMSYS Chat provide large public human--LLM interaction logs from deployed chat systems, whereas ShareChat provides a smaller publicly shared conversation corpus after English-language filtering; together they offer complementary naturalistic settings rather than experimentally matched platform conditions~\cite{zhao2024wildchat,zheng2023lmsyschat1m,yan2025sharechat}. WildChat coding contributes 31,878 conversations with 124,073 user turns and 124,623 assistant turns. WildChat writing contributes 39,534 conversations with 163,705 user turns and 164,824 assistant turns. LMSYS Chat coding contributes 32,114 conversations with 107,065 user turns and 105,174 assistant turns. LMSYS Chat writing contributes 21,023 conversations with 77,906 user turns and 76,279 assistant turns. ShareChat contributes 2,481 coding conversations with 11,541 user turns and 11,522 assistant turns and 1,539 writing conversations with 7,395 user turns and 7,363 assistant turns.

The six task settings were analysed with a common Level 1--4 annotation and analysis framework so that prevalence, contextual variation, support--engagement association and temporal ordering were computed from the same construct definitions. The manuscript reports the refreshed production outputs after the latest user-turn update pass and the subsequent LMSYS Chat and ShareChat annotation runs rather than earlier pre-update batches. Supplementary Section A.1 describes the corpus construction and refresh provenance in more detail.

Corpus filtering and classification are reported in Supplementary Table~\ref{tab:corpus_filtering_pipeline}, with retained provenance artifacts and script-hash prefixes in Supplementary Table~\ref{tab:corpus_provenance_artifacts}. In brief, coding and writing task ecologies were assigned by English-oriented LLM-assisted semantic filters rather than by native topic metadata alone. ShareChat additionally underwent a strict-English screen after task filtering, retaining only conversations whose final conversation-level language flag was English. For ShareChat, public dataset-card turns, raw CSV message rows and role-normalized analysis turns are distinct counting units; the conversion therefore uses platform-url conversations and role-labelled message rows before deriving user and assistant turns. Supplementary Table~\ref{tab:filter_classifier_validation} summarizes the available validation and audit evidence for the task, language and user-framing classifiers.

\subsection{Conversation metadata and task ecologies}

Each conversation is represented as an ordered sequence of user and assistant turns together with available metadata. The main contextual variables are data source, task domain, user framing, topic and language. Task domain and user framing are analytic variables generated by the filtering and annotation pipeline; topic and language are retained as corpus or pipeline metadata where available. Coding and writing were selected because they are common everyday LLM use cases but differ in their interactional affordances: coding often involves debugging, error diagnosis and iterative testing, whereas writing often involves drafting, revision and directive editing~\cite{keuning2018systematic,flower1981cognitive,sommers1980revision}. User framing is a conversation-level LLM-assisted label that distinguishes conversations explicitly oriented toward understanding, exploration or learning from conversations without such explicit learning-oriented framing.

Behavioural depth is measured from the number and ordering of turns in the conversation. Depth is used descriptively, as a factor of interactional opportunity and sustained exchange, not as an independent manipulation. Post-answer depth is measured as the number of turns after the first assistant response.

\subsection{Turn-level constructs}

The unit of annotation is the conversational turn. User and assistant turns are labelled with different constructs because they play different roles in the interaction. User turns are labelled for behavioural, emotional and cognitive engagement, drawing on engagement theory and the ICAP distinction between passive, active and constructive participation~\cite{fredricks2004school,chi2014icap}. The main analyses emphasize cognitive engagement and its depth levels. Passive engagement (\texttt{P}) captures reception or acknowledgement without substantial transformation. Active engagement (\texttt{A}) captures task-relevant action, follow-up or manipulation of supplied information. Constructive engagement (\texttt{C}) captures elaboration, revision, testing, extension or transformation of ideas. Constructive engagement is treated as the highest-level behavioural proxy for learning-oriented participation available in these logs, not as a direct measure of knowledge gain.

Assistant turns are labelled for scaffolding. Scaffolded support denotes assistant behaviour that structures, regulates or redirects the user's process while leaving some cognitive work with the user; the comparison category is a non-scaffolded reference response, including straightforward task fulfilment or stand-alone information provision. This definition follows the scaffolding tradition, especially contingency, transfer of responsibility and potential fading \cite{wood1976role,van2010scaffolding,quintana2018scaffolding}. Scaffolded turns can also receive intention labels (\texttt{I1} metacognitive, \texttt{I2} cognitive and \texttt{I3} affective support) and non-mutually exclusive means labels (\texttt{M1} feedback, \texttt{M2} hinting, \texttt{M3} instructing, \texttt{M4} explaining, \texttt{M5} modelling and \texttt{M6} questioning). Operational examples are provided in Supplementary Table~\ref{tab:label_examples}.

\subsection{Annotation workflow and verification}

The annotation workflow combined theory-grounded codebook development, prompt iteration, parser checks, human review and production-scale LLM-assisted labelling. We used LLM-assisted annotation as a scalable text-analysis workflow while retaining parser checks and human verification because prior work shows both the promise and the need for validation when language models are used as annotators~\cite{gilardi2023chatgpt,ziems2024can}. Prompt and post-processing rules were first developed on validation batches to reduce common failure modes, including over-labelling brief acknowledgements as active engagement and misclassifying directive assistant responses as scaffolding. After the pipeline stabilized, the same annotation logic was applied to the WildChat, LMSYS Chat and ShareChat analytic corpora. The archived annotation metadata record the provider, deployment name, API version, prompt version, retry policy, parser checks and rule-based post-processing rules used for production labelling; these details are summarized in Supplementary Section A.3. All examples reproduced in the manuscript are de-identified and paraphrased rather than copied from raw logs.

Human verification was used as a quality-control input rather than as a separate source of outcome measurement. We distinguish three validation roles. First, human--human agreement assesses whether reviewers could apply the codebook consistently. Second, production-label checks compare final or post-processed production labels with human review where a direct audit was available. Third, post-update audits check whether targeted update passes reduced known residual false positives. The final coding validation rerun combines three review rounds, with 658 task turns and 643 turns after deduplication. The final writing validation used 210 reviewed conversations. These validation sets were organized by task ecology rather than independently powered within each public corpus. Consolidated human-confirmed production-label audits covered user framing, 600 user turns and 180 assistant turns across the six corpus-by-task settings; these audits indicated acceptable or better reliability across user engagement, user framing, assistant scaffolding and support-form labels. WildChat also underwent a final targeted user-turn update check; LMSYS Chat and ShareChat were then processed with the same codebook, parser checks and post-processing logic. Agreement was summarized with per-label F1, MCC and Gwet's AC1 because the labels are sparse, prevalence-imbalanced and not mutually exclusive in all cases~\cite{matthews1975comparison,gwet2008computing}. All reported writing labels exceeded F1=0.70 in the human--human codebook check. The primary claims rely most heavily on constructive engagement, scaffolded support and support-form labels. Supplementary Tables~\ref{tab:label_validation} and \ref{tab:production_label_validation} report human--human codebook agreement and consolidated human--LLM production-label agreement, and Supplementary Section A.4 describes the final user-turn update check, the user-framing audit and corpus-transfer scope.

\subsection{Outcome construction}

The primary user-side outcome is constructive engagement. At the turn level, this is a binary indicator for whether the user turn is labelled \texttt{C}. At the conversation level, we use both the count of constructive user turns and an indicator for whether a conversation contains at least one constructive user turn. Constructive ratios are defined as constructive user turns divided by user turns.

The primary assistant-side exposure is scaffolding, operationalized as scaffolded support. At the turn level, this is the distinction between scaffolded support and a non-scaffolded reference response. At the conversation level, we use an indicator for whether a conversation contains at least one scaffolded assistant turn. Support-form analyses condition on scaffolded turns and compare conversations containing each means label with scaffolded conversations that do not contain that label. Because means labels can co-occur, these are descriptive support-form contrasts rather than mutually exclusive treatment comparisons.

\subsection{Statistical analyses}

Descriptive analyses estimate engagement prevalence, scaffolding prevalence, passive/active/constructive composition, emotional-engagement rates and behavioural-depth summaries within each task setting. Context analyses compare intentional versus unintentional conversations and coding versus writing settings. Depth analyses group conversations into behavioural-depth buckets and estimate the probability that a conversation contains at least one constructive turn. Because any-event outcomes mechanically become easier to observe in longer conversations, we also fit a grouped-binomial constructive-rate sensitivity model in which the outcome is the constructive user-turn count out of total user turns, with user framing, task ecology, length bucket and dataset fixed effects as predictors.

Conversation-level associations are analysed by comparing conversations with and without at least one scaffolded assistant turn. Constructive-ratio contrasts use turn-weighted ratios within each comparison group, defined as total constructive user turns divided by total user turns, with uncertainty from 1,000 conversation-level bootstrap resamples. The adjusted Section~2.3 count model is a setting-specific Poisson generalized linear model for the raw constructive-turn count per conversation, not for a rate. For each setting $s$, the generic mean model was $\mathbb{E}(C_i)=\exp(\alpha_s+\beta_s S_i+\gamma_s^\top X_i)$, where $C_i$ is the constructive-turn count, $S_i$ indicates at least one scaffolded assistant turn and $X_i$ contains observed conversation-level context such as user framing, conversation length, emotional-expression ratio and error or persistence markers where available. We did not use a user-turn offset in the primary count model; conversation length entered directly as a covariate. The exponentiated coefficient for scaffolded-support presence is reported as a multiplicative association with expected constructive-turn count conditional on observed covariates. As a rate sensitivity, we refitted the same setting-specific mean model with user turns represented as an exposure offset and quasi-Poisson scaled standard errors.

The companion binary model for Section~2.3 is a setting-specific logistic regression for whether a conversation contains at least one constructive user turn, using scaffolded-support presence and the same core conversation-level context available within settings. Task ecology is fixed within each setting and is therefore used in pooled/context models, not interpreted as a within-setting coefficient. Supplementary Table~\ref{tab:setting_model_covariates} lists the estimable covariate sets used for the setting-specific count, binary and offset-rate models. Model-based standard errors and 95\% confidence intervals are shown for the adjusted Poisson and logistic estimates. We checked Poisson overdispersion using the Pearson $\chi^2$ statistic divided by residual degrees of freedom; the dispersion estimates ranged from 1.16 to 2.06 across the six settings. A quasi-Poisson sensitivity that scaled the Poisson standard errors by this dispersion preserved the direction and statistical distinguishability of the scaffolded-support association in every setting, and the offset-rate sensitivity likewise preserved positive scaffolded-support associations in every setting. Negative-binomial models were not used as the primary specification because the estimand was the common pre-specified multiplicative association under the same mean model across settings; the overdispersion and offset checks were used to assess variance and exposure-specification sensitivity rather than to change the main estimand. Binary outcomes in other analyses, including whether the next user turn is constructive and whether the next assistant turn is scaffolded, are modelled with logistic regression or reported as conditional probability contrasts, depending on the estimand~\cite{mccullagh1989generalized,agresti2013categorical}. Percentage-point contrasts are reported as effect sizes; adjacent-turn contrasts use 1,000 conversation-cluster bootstrap resamples because multiple turn pairs can come from the same conversation. Support-form bracket tests in Fig.~\ref{fig:support_form_supply} use Benjamini--Hochberg false-discovery-rate adjustment within each displayed family of six support-form comparisons. Supplementary Tables~\ref{tab:model_specs}--\ref{tab:temporal_estimands} summarize the main model families and temporal estimands.

To assess whether the adjacent-turn pattern was reducible to a single feature family, we also fitted integrated logistic regressions that combined user-state features, assistant-scaffolding features and adjacent-turn context in one model. The outcome was whether the next user turn was constructive. The pooled model included dataset fixed effects; a sensitivity replaced these with recoverable assistant model or source-family fixed effects, using exact \texttt{chat\_model} labels where available, LMSYS Chat model names recovered from conversation identifiers and ShareChat assistant/source families recovered from conversation identifiers. Standard errors for these adjacent-turn regressions were clustered by conversation. We report broad scaffolded-support models separately from support-form decompositions because M1--M6 are nested descriptors of scaffolded support rather than independent exposures. Nested likelihood-ratio block tests first add broad scaffolded-support presence and then test whether support-form descriptors add signal within scaffolded support. These model/source checks are interpreted as robustness analyses rather than model-causal comparisons.

Temporal analyses focus on adjacent-turn ordering. The primary analyses compare the probability that the next user turn is constructive after scaffolded versus non-scaffolded assistant turns, estimate the same contrast within prior user states and estimate the probability that the next assistant turn is scaffolded after different prior user states. We also fit adjacent-turn regressions with prior-state $\times$ support-form interactions to test whether coupling depends jointly on the user's immediately preceding engagement state and the form of scaffolding provided. Standard errors for these adjacent-turn regressions were clustered by conversation. Coarser within-conversation before--after summaries around the first scaffolded assistant turn are reported as aggregate temporal diagnostics. They are used to contextualize the adjacent-turn claim rather than to estimate cumulative learning trajectories, because the first scaffolded turn may be selected by task difficulty, user motivation or an already developing exchange.


% Coarser within-conversation before--after summaries are reported only as exploratory boundary checks, not as evidence of cumulative learning trajectories, because the first scaffolded turn may be selected by task difficulty, user motivation or an already developing exchange.



\subsection{Robustness, privacy and LLM use}

WildChat robustness analyses repeat the main patterns across user-facing model snapshots and coarse model families. These checks assess whether the headline patterns are reducible to a single WildChat model label. They are not interpreted as model-causal comparisons, because model labels are not experimentally assigned and may be confounded with time, user population, task mix and logging context.

The analyses are bounded by the observational design and by the nature of the labels. They support claims that everyday human--LLM conversations contain measurable learning-oriented engagement; that constructive engagement varies by intent, task ecology and depth; that assistant scaffolding is associated with richer constructive participation; and that support and engagement show adjacent-turn reciprocal ordering. They do not support claims that scaffolded support causes durable learning gains, that all constructive engagement reflects actual knowledge change or that the labelled conversational states perfectly measure learning itself.

LLMs were used as objects of study and as annotation tools. Production labels were generated with LLM-assisted workflows and then checked through parser rules, targeted update scripts and human validation. The authors also used LLM-based tools to assist with code drafting, debugging, analysis workflow organization and language editing. LLM-generated assistance was not used as a substitute for human verification of the validation samples or for interpretation of the statistical results.

All reported results are aggregate turn-level or conversation-level analyses. The study does not link conversations into user histories, estimate user-level longitudinal trajectories or publish raw conversation text; temporal analyses are confined to ordering within a conversation.
